\documentclass[10pt,a4paper]{amsart}
\usepackage[margin=19mm]{geometry}
\usepackage{amsmath,amssymb,mathtools}
\usepackage[scr=rsfs,cal=euler]{mathalfa}
\usepackage{xfrac}
\usepackage[unicode,hidelinks,bookmarks=false]{hyperref}
\newcommand{\ie}{i.e.\ }
\newcommand{\cf}{cf.\ }
\newcommand{\resp}{respectively\ }
\newcommand{\Subsection}{\subsection}
\providecommand{\nobreakdash}{\nobreak\hbox{-}}
\setlength{\emergencystretch}{2em}
\multlinegap0pt
\usepackage{enumitem}
\usepackage{mathrsfs}
\theoremstyle{plain}
\newtheorem{theorem}{Theorem}[section]
\newtheorem{proposition}[theorem]{Proposition}
\newtheorem{lemma}[theorem]{Lemma}
\newtheorem{corollary}[theorem]{Corollary}
\newtheorem{definition}[theorem]{Definition}
\newtheorem{example}[theorem]{Example}
\theoremstyle{remark}
\newtheorem{remark}[theorem]{Remark}
\newcommand{\la}{\label}
\newcommand{\fr}{\frac}
\newcommand{\na}{\nabla}
\newcommand{\be}{\begin{equation}}
\newcommand{\ee}{\end{equation}}
\newcommand{\Rr}{{\mathbb R}}
\title{On putative self-similarity for incompressible 3D Euler}
\author{Peter Constantin, Mihaela Ignatova and Vlad Vicol}
\date{Source: arXiv:2602.17570v3 (CC BY 4.0); the authors' own native \LaTeX{} source, set here in the editorial AMS \texttt{amsart} wrapper (the source's own class is \texttt{amsproc}, which is not redistributed).}
\begin{document}
\maketitle
\noindent\textit{Editorial scope.} The counted claim of this unit is the article's Proposition 3.1 (source label \texttt{nse}), the viscous statement that a locally self-similar inner profile with exponent $\gamma$ can only blow up if $\gamma\leq 1/2$, delivered together with the authors' complete original proof. Retained uncounted support, in source order, is the whole source-local core that the counted claim needs: the passage of Section~3 that precedes it, including the equation and assumption labels that the statement and the proof refer to (source0.tex lines 266--377), and the remark that discusses the hypotheses of the proposition and is placed between the statement and its proof (source lines 399--430). Every retained passage is a verbatim copy of the authors' own source lines and every cross-reference inside them resolves inside the delivered document. Printed numbering: the authors' own theorem-environment declarations are carried unchanged, with the proposition sharing the counter of the theorem environment and being numbered by section as in the authors' preamble, and the counters are advanced so that the delivered PDF prints the counted claim as Proposition~3.1, exactly as the published article does. Omitted: the title block, abstract and introduction of the article, the parts of the text that precede line 269, and everything after the counted proof. Nothing of the counted statement or of the counted proof is omitted. Class substitution: the source class is the AMS \texttt{amsproc} class with the authors' own packages; the wrapper is AMS \texttt{amsart} carrying the authors' own macro and theorem declarations as well as the \texttt{enumitem} package used by the authors' itemized remark. Reference handling: the two works cited by the retained passages are transcribed verbatim from the article's own bibliography, which is part of the authors' own TeX and is bundled unchanged inside the source file, and they keep the printed numbers 17 and 22 that the citing sentences carry in the published article. No mathematics is written, repaired or re-derived here.

\setcounter{section}{3}
\setcounter{theorem}{0}
\begin{equation}
\partial_t u + u \cdot \nabla u+ \nabla p = 0 , \qquad 
\nabla \cdot u =0  , 
\label{eq:Euler:velocity}
\end{equation}
admit a globally self-similar singularity at a time $ T_* > 0$, at a single point $x_* \in \mathbb{R}^3$. 

\subsection{The self-similar ansatz}
\label{sec:standing:SS:ass}
Due to Galilean symmetry, we assume without loss of generality that the singularity occurs at the space location $x_*=0$, and 
that the similarity profile for the velocity vanishes at this point.  Due to time-rescaling symmetry, we  take without loss of generality $T_*=1$. 
That is, we investigate the globally self-similar ansatz:
\begin{equation}
\begin{aligned}
&u(x,t) = (1-t)^{\gamma-1} U(y)
, 
\qquad
p(x,t) = (1-t)^{2(\gamma-1)} P(y) 
,
\\
&\omega(x,t) = \frac{1}{1-t} \Omega(y) 
, 
\qquad
y = \frac{x}{(1-t)^\gamma}
.
\end{aligned}
\label{eq:SS:ansatz}
\end{equation}
This ansatz is less restrictive than it seems. A more general form, $\omega (x,t) =
A(t)\widetilde \Omega\left(\fr{x-x(t)}{\ell(t)}, \tau (t)\right)$ leads to the form above. Indeed, writing the Euler equations in vorticity form and using the fact that the velocity and vorticity are related by the linear relation
$\omega = \na\times u$, we see that the nonlinear term $u\cdot\na\omega-\omega\cdot\na u$ scales with amplitude $A^2$ under the general ansatz above. We thus arrive at the ODE $\dot A = A^2$ where $\dot A$ is a time derivative {\footnote{ This time derivative can be with respect to an arbitrary internal clock, that is: $\fr{d}{dt}$ can be replaced by $\fr{1}{f(t)}\fr{d}{dt}$ where $f(t)$ is an arbitrary non-vanishing function.}}. This leads inevitably to $A(t) = (T-t)^{-1}$. The rest follows dividing the Euler equation by $A^2$ and taking the simplest situation, when the ensuing coefficients are constant. Then the constant coefficient in front of $y\cdot\na \Omega$ leads to $\ell(t) = \ell_0\left(1 -\fr{t}{T}\right)^{\gamma}$. The term due to $x(t)$ is integrated, because $A$ and $\ell$ are known. It follows that $x(t) = x_* - \fr{c}{\gamma}\ell(t)$ where $c$ is a constant, and this is dealt with by translating $\widetilde \Omega$.

 
\subsubsection{The stationary self-similar PDE}
The similarity profiles $U,P,\Omega$ appearing in~\eqref{eq:SS:ansatz} satisfy the (stationary) self-similar Euler equation in velocity form
\begin{equation}
(1-\gamma) U + \gamma (y \cdot \nabla) U + (U \cdot \nabla) U + \nabla P = 0,
\qquad 
\nabla \cdot U = 0 ,
\label{eq:profile:U}
\end{equation}
or equivalently, the (stationary) self-similar Euler equation in vorticity form
\begin{equation}
\Omega + \gamma (y \cdot \nabla) \Omega + (U \cdot \nabla) \Omega = (\Omega \cdot \nabla) U
,
\qquad 
\Omega = \nabla \times U
,
\qquad 
\nabla \cdot U = 0 .
\label{eq:profile:Omega}
\end{equation}
As mentioned earlier, due to Galilean symmetry (translation invariance), we assume throughout this section that the self-similar velocity profile $U$ vanishes at the origin. We also only consider profiles $U$ which are sub-linear at infinity.\footnote{
These are the profiles for which we may (typically) employ a localization or cutoff  procedure at
large values of $y$, to ensure that the velocity field $u$ in original $(x,t)$ coordinates has finite kinetic energy.}
We summarize these properties as
\begin{equation}
U(0) = 0,
\qquad
\lim_{|y|\to \infty} |y|^{-1} |U(y)| = 0 \,.
\label{eq:U:non:negotiable} 
\end{equation}

\subsubsection{Space rescaling and normalization}
Note that if $U$ and $\Omega$ are a solution of \eqref{eq:profile:Omega} on $\mathbb{R}^3$, then   
$U_{\lambda}(y)= \lambda U(\tfrac{y}{\lambda})$ and $\Omega_{\lambda}(y) =  \Omega(\tfrac{y}{\lambda})$, 
are also a solution of \eqref{eq:profile:Omega} on $\mathbb{R}^3$.
Under this rescaling, we have\footnote{As we shall see later (cf.~\eqref{eq:ass:medium:rare}), $\Omega$ is only expected to lie in $L^p(\mathbb{R}^3)$ when $p > 3 \gamma$.} 
\begin{equation}
\|\nabla^2 U_\lambda\|_{L^\infty(\mathbb{R}^3)}
= \tfrac{1}{\lambda} \|\nabla^2 U\|_{L^\infty(\mathbb{R}^3)}
,\quad
\|\nabla \Omega_\lambda\|_{L^\infty(\mathbb{R}^3)}    
= \tfrac{1}{\lambda}  \|\nabla \Omega\|_{L^\infty(\mathbb{R}^3)}    
,\quad 
\| \Omega_\lambda\|_{L^p(\mathbb{R}^3)}    
=  \lambda^{\frac{3}{p}} \|\Omega\|_{L^p(\mathbb{R}^3)}   
.
\label{eq:Omega:scaling}
\end{equation}
In order to fix units, throughout this section we consider profiles which are normalized via~\eqref{eq:Omega:scaling} to satisfy
\begin{equation}
\| \nabla^2 U \|_{L^\infty(\mathbb{R}^3)}  = 1 .
\label{eq:U:scaling} 
\end{equation}

 

\subsubsection{Behavior of the profiles at space infinity}
While the assumption of  sublinearity of $U$ at infinity (cf.~\eqref{eq:U:non:negotiable}) is necessary for any reasonable self-similar profile, the self-similar PDEs~\eqref{eq:profile:U} and~\eqref{eq:profile:Omega} impose more stringent assumptions. Indeed, \eqref{eq:U:non:negotiable} dictates that the term $\gamma y \cdot \nabla$ is stronger than the term $U \cdot \nabla$ for large values of $|y|$. Thus, the behavior of $U$ and $\Omega$ as $|y|\to \infty$ must respect the kernel of the operators $(1-\gamma) + \gamma y \cdot \nabla$ (for $U$) and $1 +\gamma y \cdot \nabla$ (for $\Omega$). This formally suggests that the leading order behavior of the self-similar profiles in the far field is given by 
\[
|U| \sim |y|^{\frac{\gamma-1}{\gamma}} ,
\qquad\mbox{and}\qquad
|\Omega| \sim |y|^{-\frac{1}{\gamma}},
\qquad\mbox{as} \qquad |y|\to \infty.
\]
It is convenient to quantify the above asymptotic descriptions. Recalling that $U(0) = 0$, we assume that there exists a constant $\mathsf{C}_\flat > 0 $ such that
\begin{equation}
  |U(y)| \leq \mathsf{C}_\flat |y| \langle y\rangle^{-\frac{1}{\gamma}},
  \qquad
  \mbox{and}
  \qquad
  |\Omega(y)| + |\nabla U(y)| \leq \mathsf{C}_\flat  \langle y\rangle^{-\frac{1}{\gamma}},
  \qquad
  \mbox{for all}
  \qquad y\in\mathbb{R}^3.
  \label{eq:ass:medium:rare}
\end{equation}

\subsection{A remark for the viscous case.}
We consider here hypothetical  blow up  for the 3D incompressible Navier-Stokes equation {(set $\mathsf{RHS}_{\eqref{eq:Euler:velocity}} = \Delta u$ instead of $0$)}, where the vorticity $\omega(x,t)$ is approximately self-similar.


\begin{proposition}\label{nse}
Let $\omega(x,t)$ be a solution of the 3D incompressible Navier-Stokes equation which can be written as
\begin{equation*}
\omega(x,t) 
= {\omega_{\rm in}(x,t)} + \omega_{\rm out}(x,t),
\end{equation*}
where $\omega_{\rm out}$ is regular, in the sense that
\begin{equation}
\int_0^{{1}} \|\omega_{\rm out}(\cdot,t)\|_{L^2(\Rr^3)}^4dt 
\le \Gamma<\infty,
\label{omegaoutfine}
\end{equation}
and where $\omega_{\rm in}$ is supported in $B_{\ell(t)}(x(t))$
 with $\ell(t)\sim (1-t)^{\gamma}$, {and satisfies}
\begin{equation}
\|\omega_{\rm in}(\cdot,t)\|_{L^2(B_{\ell(t)}(x(t)))} \le  C (1-t)^{-1+\fr{3\gamma}{2}}
\label{omegainbad}
\end{equation}
{for some $C<\infty$.}
Then, if the solution blows up at time $t=1$, it follows that $\gamma\leq \frac{1}{2}$.
\end{proposition}


\begin{remark}
The assumption \eqref{omegainbad} is satisfied if
\[
\omega_{\rm in}(x,t) \sim
\frac{1}{{1}-t}\Omega\left(\fr{x-x(t)}{({1}-t)^{\gamma}}\right) 
\qquad\mbox{in} \qquad B_{\ell(t)}(x(t)), 
\qquad\mbox{as} \qquad t \to 1^-,
\] 
with $|\Omega(y)| \le  C$ for $|y|\le 1$, {and $\ell(t) \sim (1-t)^\gamma$}. 
The assumption \eqref{omegaoutfine} is satisfied {for instance} if 
$\omega_{\rm out}(x,t) = 0$ for $|x-x(t)|\le \ell(t)$,
\[
|\omega_{\rm out}(x,t)| \le C_1 |x-x(t)|^{-\fr{1}{\gamma}}, 
\qquad \mbox{for} \qquad \ell(t) \leq |x-x(t)|\leq 1,
\]
or, more generally, if
\be
\int_0^1\|\omega_{\rm out}(t)\|^4_{L^2(B_1(x(t)))}dt 
\leq C_2 <\infty,
\la{outin}
\ee
and if, for instance
\[
\sup_{|x - x(t)|\ge 1}|\omega_{\rm out}(x,t)| \le C_3.
\] 
Indeed, because $\omega_{\rm out}(x,t) = \omega(x,t)$ for all $|x-x(t)|\ge 1$ we have  $|\omega_{\rm out}(x,t)|^2 \le C_3|\omega(x,t)|$ for 
$|x-x(t)|\ge 1$. Then $\|\omega_{\rm out}(t)\|_{L^2(\Rr^3\setminus B(x(t),1))}^2 \le C_3\|\omega(t)\|_{L^1(\Rr^3)}$. As it is well known, $\sup_t \|\omega (t)\|_{L^1(\Rr^3)}$ is bounded in terms of the initial data~\cite{ConstantinArea}, so the assumptions imply that $\omega_{\rm out}\in L^4(dt; L^2(\Rr^3))$. 

Thus, the assumption~\eqref{omegaoutfine}  is satisfied if $\omega$ is a smooth vorticity matched to an inner self-similar blow up ansatz with asymptotic behavior at infinity of the type~\eqref{eq:ass:medium:rare}. The main point is that both the inner blow up profile and the matching vorticity magnitude majorant  $z(x,t) = |x-x(t)|^{-\fr{1}{\gamma}}\mathbf 1_{\{\ell(t)\le |x-x(t)| \le 1\}}$ have the property  that  they belong to $ L^4(dt; L^2(\Rr^3))$ if $\gamma>\fr{1}{2}$, which is a class of non-blow up.
\end{remark}


\begin{proof}[{Proof of Proposition~\ref{nse}}]
We add the contributions {from the inner~\eqref{omegainbad} and outer~\eqref{omegaoutfine} pieces} and deduce
\begin{equation*}
\int_0^{{1}} \|\omega (\cdot,t)\|_{L^2(\Rr^3)}^4dt \le 16\left(\Gamma + \fr{C^4}{6\gamma-3}\right).
\end{equation*}
If $\gamma>\fr{1}{2}$, the right hand side is finite, and there is no blow up because of a well-known regularity criterion~\cite{CFbook}.
\end{proof}

\begin{thebibliography}{99}
\bibitem[17]{ConstantinArea} P.~Constantin. \newblock Navier-Stokes equations and area of interfaces. \newblock \textit{Communications in Mathematical Physics}~\textbf{129} (2):241--266, 1990.
\bibitem[22]{CFbook} P.~Constantin, C.~Foias. \newblock {Navier-Stokes Equations}. \newblock University of Chicago Press, Chicago, 1988.
\end{thebibliography}
\end{document}
